% !PRINTED 419
\weberpage{419}
%<<CONT
particular values of the argument $\omega$ that we considered in the preceding section. On these depend the moduli of the elliptic functions which admit a complex multiplication. After \gesperrt{Kronecker} they are called \gesperrt{singular moduli}. Accordingly we shall speak also of the \gesperrt{singular values of the modular functions}, and in particular of the \gesperrt{singular invariants} $j(\omega)$, and we understand by these the values that these functions take when $\omega$ is the \gesperrt{root} (with positive imaginary part) \gesperrt{of a quadratic form with negative discriminant}.

If $\omega$ satisfies equation (3), § 114, and $a, b, c, \partial, n$ are determined as in the preceding section, then there follows first of all:

\weq{1}{j(\omega) = j\left(\frac{c + \partial\,\omega}{a + b\,\omega}\right),}

and if therefore

\weq{2}{F_n(u, v) = 0}

is the invariant equation belonging to the degree of transformation $n$ (§ 69), then (2) is satisfied if $a, b, c, \partial$ are without a common divisor and

\weq{3}{u = j(\omega), \quad v = j\left(\frac{c + \partial\,\omega}{a + b\,\omega}\right) = u}

is set. Thereupon we arrive at the first principal theorem of this theory:

III. \gesperrt{The singular value}
\[
u = j(\omega)
\]
\gesperrt{is a root of the algebraic equation:}

\weq{4}{F_n(u, u) = 0.}

If conversely $u$ is a root of equation (4), then, when $\omega$ is determined from the equation $j(\omega) = u$, one of the values of $v$ that satisfy equation (2) is likewise $= u$, and there subsists therefore an equation of the form (1); by § 53 this has the consequence that $\omega$ is equivalent to $\dfrac{c + \partial\,\omega}{a + b\,\omega}$, and from this there results an equation of the form § 114, (3). Thus we have the complement to III.:

IV. \gesperrt{Every root of (4) is a singular invariant.}
